\section*{Chapter \ref{ch:rc:operational-risk-and-rogue-trading} --- Operational Risk and Rogue Trading}

\begin{solution}{exo:rc:operational-risk-and-rogue-trading:1}
$0.12\times1+0.15\times9 = \text{EUR}~1.47$ billion.
\end{solution}

\begin{solution}{exo:rc:operational-risk-and-rogue-trading:2}
Then $(\mathrm{LC}/\mathrm{BIC})^{0.8}=1$ and $\mathrm{ILM}=\ln(e-1+1)=\ln e=1$: a bank whose losses are in line with its
size holds the BIC. More losses raise the multiplier, fewer lower it, both less than proportionally.
\end{solution}

\begin{solution}{exo:rc:operational-risk-and-rogue-trading:3}
The deferred-settlement rule (settlement dates far beyond the market's two days), together with the late
booking rule, since the FSA found both in the scheme.
\end{solution}

\begin{solution}{exo:rc:operational-risk-and-rogue-trading:4}
Collateral of a tenth of the money lent means either the clients were given unsecured credit no one
approved, or the money was not going to clients at all. Either way funding requests that cannot be
reconciled with known positions and clients should stop the funding until they are explained.
\end{solution}

\begin{solution}{exo:rc:operational-risk-and-rogue-trading:5}
The rules flag genuine trades for unrelated reasons (a late booking here, an off-market fill there), so
two rules rarely coincide on a genuine trade; a concealment scheme breaks several rules at once by design.
On the blotter, false alarms fall from 7.1\% to 0.05\% while hits fall only from 88.3\% to 61.7\%.
\end{solution}

\begin{solution}{exo:rc:operational-risk-and-rogue-trading:6}
A concealed position needs its \omterm{def:rc:operational-risk-and-rogue-trading:rogue}{fictitious trades} rolled, amended and explained every day. With the trader
away and locked out, the positions are run by someone else, and breaks, margin calls and unmatched trades
surface.
\end{solution}

\begin{solution}{exo:rc:operational-risk-and-rogue-trading:7}
\texttt{audit\_demo()} returns \texttt{True} before and \texttt{False} after the edit. Rewriting the chain
from the edited entry onward produces a consistent chain with a different latest hash; if that hash was
copied elsewhere (another system, a daily report, a third party), the rewrite is detected by comparison.
\end{solution}

\begin{solution}{exo:rc:operational-risk-and-rogue-trading:8}
Steady, high profits with low measured risk are what a hidden position or fictitious hedge looks like:
the returns are too good for the risk the firm sees. The right response is to explain the P\&L by strategy and risk before raising any limit.
\end{solution}

\begin{solution}{pb:rc:operational-risk-and-rogue-trading:1}
\textbf{1.} $6.4/49 = 13.1\%$ (the net 4.9 billion mixes in the 2007 gains).
\textbf{2.} $6.4/0.05 = \text{EUR}~128$ billion.
\textbf{3.} Selling tens of billions of shares in a few days is a large share of the market's volume; prices
moved against the bank under its own selling.
\textbf{4.} $49\times0.25/\sqrt{250}\times2.326 = \text{EUR}~1.80$ billion.
\textbf{5.} The \omterm{def:rc:operational-risk-and-rogue-trading:rogue}{fictitious trades} offset the real ones in the firm's records, so the net position the VaR
saw was small.
\textbf{6.} Trades in the opposite direction with internal or fictitious counterparties, or forward trades
that would never settle.
\textbf{7.} Internal unmatched, deferred settlement, cancel-and-amend near reporting, and late booking.
\textbf{8.} Nothing, as the trades were cancelled before their confirmation date; confirming at booking would have shown trades no counterparty recognised.
\textbf{9.} Margin and cash flows on the real positions, far larger than the net position explained.
\textbf{10.} Without the trader to roll and amend the \omterm{def:rc:operational-risk-and-rogue-trading:rogue}{fictitious trades}, they break.
\textbf{11.} 88.3\% hit and 7.1\% false alarms with any rule; 61.7\% and 0.05\% with two rules.
\textbf{12.} About 200 a year with one rule (142 false, 53 real), about 40 with two, for a blotter of
2\,060 trades; control functions independent of the desk investigate them.
\textbf{13.} Too many false alarms, closure on the explanation of the person concerned without evidence, and
no escalation of repeats.
\textbf{14.} It shows every booking, amendment and closure with its time and author, and proves the record
was not edited.
\textbf{15.} Margin against positions, exchange position reports against the book, and gross against net.
\textbf{16.} A first loss, hidden to avoid its consequences, grows as the trader bets to recover it.
\textbf{17.} A rule on cancelled trades with distant value dates, at booking: confirmation came too late.
\textbf{18.} Through the loss component: the losses enter the average for ten years and raise the ILM.
\textbf{19.} Named result: \emph{the January 2008 arithmetic}: an unwinding loss of EUR~6.4 billion on
49 billion is an average adverse move of 13.1\%; the first control in the chain that would have caught it is at
booking: flagging trades with distant value dates that are cancelled before confirmation.
\textbf{20.} It is followed through: every alert closed with evidence, independently of the people it controls.
\end{solution}

\begin{solution}{iq:rc:operational-risk-and-rogue-trading:1}
The risk of loss from failed processes, people, systems and external events, including legal risk.
Market and credit risk are taken deliberately for a return; \omterm{def:rc:operational-risk-and-rogue-trading:oprisk}{operational risk} comes with doing business,
has no upside, and is heavy-tailed.
\iqlookfor{the definition and the difference in nature.}
\end{solution}

\begin{solution}{iq:rc:operational-risk-and-rogue-trading:2}
\omterm{def:rc:operational-risk-and-rogue-trading:sod}{Segregation of duties}, independent booking and confirmation, reconciliation of cash, margin and
positions, P\&L explain, limits, \omterm{def:rc:operational-risk-and-rogue-trading:sod}{block leave}, and investigated alerts. At Barings the trader ran the back
office, funding was not reconciled with positions, and exceptional profits were not questioned.
\iqlookfor{the controls and the documented failures.}
\end{solution}

\begin{solution}{iq:rc:operational-risk-and-rogue-trading:3}
Feeds from booking, confirmation, settlement and market data; rules like the chapter's plus behaviour
baselines per trader; scores and thresholds tuned on false alarms; case management with evidence and
escalation; an append-only audit log; independent ownership.
\iqlookfor{data, rules, scoring, investigation and independence.}
\end{solution}

\begin{solution}{iq:rc:operational-risk-and-rogue-trading:4}
Do not ignore it; ask the desk head and risk for the explanation by strategy, and escalate to control
functions if none comes. Unexplained, steady profit is a warning sign.
\iqlookfor{escalation, not investigation on one's own.}
\end{solution}

\begin{solution}{iq:rc:operational-risk-and-rogue-trading:5}
Capital is the BIC (marginal coefficients of the business indicator) times the ILM from ten years of
losses. Weaknesses: size is a crude proxy of risk, past losses predict tail events poorly, and it gives
little incentive for better controls.
\iqlookfor{the formula and its limits.}
\end{solution}

\begin{solution}{iq:rc:operational-risk-and-rogue-trading:6}
Too many false alarms, closure on the trader's explanation, no escalation of repeats, and investigators
who depend on the desk. Fix: combined scores, evidence required to close, repeat tracking, and
independent investigators with standing.
\iqlookfor{alert fatigue and independence.}
\end{solution}
